

%قبل از شروع مسئله ترجیحاً چیزی ننویسید. اما اگر خواستید بنویسید، در قالب لیستی (فقرات و شمارش) بنویسید.
%برای اضافه کردن نمره، بعدی از شروع‌مسئله، [x] را اضافه کنید و به جای x  نمره مورد نظر را قرار دهید. برای زیرمسئله‌ها هم به همین صورت است.
%\شروع‌مسئله[10]
\شروع‌مسئله []
%\متن‌سیاه{سوال عملی اول:}\\
%\متن‌سیاه{دراپ‌تست}\\

مسیح یک رشته s از اعداد ۰ و ۱ دارد. او می‌داند به یک زیررشته جادویی گفته می‌شود اگر عدد دودویی که آن زیررشته نشان می‌دهد بر ۷ بخش‌پذیر باشد. حال او می‌خواهد بداند کدام یک از زیررشته‌های رشته‌اش جادویی‌اند یا با تغییر بعضی کاراکترهای رشته جادویی می‌شوند؛ برای همین در هر مرحله یکی از پرسش‌های زیر را از شما می‌کند: \\
۱. \متن‌لاتین{0 l r} : باقی‌مانده‌ی عدد دودویی که زیررشته‌ی l تا r رشته‌اش یعنی \متن‌لاتین{$s_l$$s_{l+1}$...$s_r$} ‌ نشان می‌دهد
 را بر عدد ۷ به او بگویید تا بداند زیررشته‌اش جادویی است یا نه.\\
۲. \متن‌لاتین{1 p v } : مقدار خانه‌ی p رشته‌ی او را به مقدار v تغییر دهید.\\
توجه کنید که اندیس‌گذاری رشته از صفر شروع می‌شود. 

\متن‌سیاه{ورودی:}\\
هر ورودی شامل چند تست است. در خط اول t که تعداد تست‌ها می‌آید.\\
در خط اول هر تست رشته‌ی دودویی s به طول n شما داده می‌شود. در خط بعد  \متن‌لاتین{q} آمده که تعداد درخواست‌ها را مشخص می‌کند.\\ 
در هر یک از q خط بعدی سه عدد آمده که مطابق یکی از درخواست‌های گفته شده است.\\


\متن‌سیاه{خروجی:}\\
به ازای هر کدام از درخواست‌های نوع اول مقدار خواسته شده را چاپ کنید.\\

\متن‌سیاه{محدودیت‌ها:}\\
$1 \leq t \leq 100$\\
$1 \leq n, q \leq 10^{4}$\\
$0 \leq l \leq r \leq n - 1$\\
$0 \leq p \leq n - 1$\\
$0 \leq v \leq 1$\\

\متن‌سیاه{ورودی و خروجی نمونه:}\\

\begin{latin}
	\begin{tabularx}{\textwidth}{|X | X|}
		\hline
		\textbf{stdin} &
		\textbf{stdout}\\
		\hline
		1 & 5\\
		01010101 & 6\\
		5 & 5\\
		0 1 3 &\\
		1 2 1 & \\
		0 2 5&\\
		1 3 1&\\
		0 0 7&\\
		\hline
	\end{tabularx}
\end{latin}

%}
% صورت اصلی مسئله را در این‌جا بنویسید.
%اگر مسئله چند قسمت داشت به این صورت عمل کنید:
%\شروع{زیرمسائل}
%\زیرمسئله[0] % مسئلۀ ۱آ
%\زیرمسئله % مسئلۀ ۱ب
%\زیرمسئله % مسئلۀ ۱پ
%\زیرمسئله % مسئلۀ ۱ت
%\پایان{زیرمسائل}



%%%%%%%%%%%%%%%%%%%%%%%%%%%%%%%%%%%%%%%%%%%%%%%%%%%%%
% برخی دستورات مفید برای تمرین‌ها. %
%%%%%%%%%%%%%%%%%%%%%%%%%%%%%%%%%%%%%%%%%%%%%%%%%%%%%

% فرمول‌های نمونه:
%$f(x) = \Theta(x)$
%$T(x, y) \leq \log(x) + 2^y + \binom{2n}{n}$
% {\tt code\_function}


% لیست موارد بدون شماره‌گذاری:
%\شروع{فقرات}
%    \فقره مورد اول 
%        \شروع{فقرات}
%            \فقره مورد اول داخلی 
%                \شروع{فقرات}
%                    \فقره اولین مورد خیلی داخلی 
%                    \فقره دومین مورد خیلی داخلی 
%                \پایان{فقرات}
%            \فقره مورد دوم داخلی 
%        \پایان{فقرات}
%    \فقره مورد دوم 
%\پایان{فقرات}

% لیست موارد شماره‌گذاری شده:
%\شروع{شمارش}
%    \فقره یک مورد 
%    \فقره یک مورد دیگر 
%    \فقره باز هم یک مورد دیگر
%\پایان{شمارش}

% برای نوشتن یک فرمول چند خطی که یک شماره بخوره. برای حالت‌های دیگر به سایت زیر مراجعه کنید:
% https://www.overleaf.com/learn/latex/Aligning_equations_with_amsmath
%\begin{align} 
%    \begin{split}
%        (x+y)^3 &= (x+y)^2(x+y) \\
%                &= (x^2+2xy+y^2)(x+y) \\
%                &= (x^3+2x^2y+xy^2) + (x^2y+2xy^2+y^3) \\
%                &= x^3+3x^2y+3xy^2+y^3
%    \end{split}                                 
%\end{align}

% نوشتن ماتریس و آرایه دوبعدی
%\begin{align}
%    A = 
%    \begin{bmatrix}
%        A_{11} & A_{21} \\
%        A_{21} & A_{22}
%    \end{bmatrix}
%\end{align}

% اضافه کردن تصویر با شرح. برای اضافه کردن تصویر بدون شرح، فقط دستور درج‌تصویر را استفاده کنید.
%\شروع{شکل}[h!]
%\تنظیم‌ازوسط
%  \درج‌تصویر[پهنا=0.5\پهنای‌متن]{img.jpg}
%\شرح{شرح تصویر}
%\برچسب{تصویر:۱}
%\پایان{شکل}
%\پایان{زیرمسائل}

%کد زیر را ببینید.
%%\lstinputlisting{source_filename.py}
%\begin{latin}
%\begin{lstlisting}
%# Import the modules
%import sys
%import random
%
%ans = True
%
%while ans:
%question = raw_input("Ask the magic 8 ball a question: (press enter to quit) ")
%
%answers = random.randint(1,8)
%
%if question == "":
%sys.exit()
%
%elif answers == 1:
%print "It is certain"
%
%elif answers == 2:
%print "Outlook good"
%
%elif answers == 3:
%print "You may rely on it"
%
%elif answers == 4:
%print "Ask again later"
%
%elif answers == 5:
%print "Concentrate and ask again"
%
%elif answers == 6:
%print "Reply hazy, try again"
%
%elif answers == 7:
%print "My reply is no"
%
%elif answers == 8:
%print "My sources say no"
%\end{lstlisting}
%\end{latin}

	